\pojem{Navadna diferencialna enačba $k$-tega reda} je enačba oblike
\[
  \vec{F}(t, x(t), x'(t), \ldots, x^{(k)}(t)) = 0,
\]
kjer je $x: \R^n \to \R^n$ funkcija, ki reši enačbo.
Če je $x$ vektorska funkcija ($n > 1$), enačbi pravimo \pojem{sistem NDE}.
Za enačbo pravimo, da je \pojem{avtonomna}, če $F$ ni eksplicitno odvisna od
$t$.
Neavtonomen sistem lahko spravimo v avtonomnega tako, da uvedemo še eno odvisno
spremenljivko $v(t)$, ter enačbo $\dot{v} = 1$.

\section{Prvi integral}%

Splošna rešitev enačbe $y' = f(x, y)$ je enoparametrična družina funkcij $y =
y(x, C)$.
Če obstaja funkcija $u(x, y)$, za katero velja $u(x, y(x, C)) = \text{konst.}$
za vsak $C$, jo imenujemo \pojem{prvi integral enačbe}.
Vsaka implicitno podana krivulja $y(x)$ z enačbo $u(x, y) = D$, kjer je $u$ prvi
integral neke diferencialne enačbe, je rešitev te diferencialne enačbe.
Iskanje prvega integrala je torej kvečjemu močnejše od reševanja diferencialne
enačbe.

Če imamo dano vektorsko polje $F = (P, Q): \R^2 \to \R^2$, lahko poiščemo
ortogonalne krivulje nanj z diferencialno enačbo
\[
  y' = - \frac{P}{Q}.
\]
Če je polje potencialno s potencialom $u$, iste krivulje poiščemo z
\[
  y' = - \frac{\partial_x u}{\partial_y u},
\]
torej je $u$ prvi integral te enačbe.
Ker potencial načeloma znamo poiskati z integralom, lahko enačbo rešimo.
Če pa polje ni potencialno, mora obstajati funkcija $\lambda = \lambda(x, y)$,
ki jo imenujemo \pojem{integrirujoči množitelj}, za katero je polje $(\lambda P,
\lambda Q)$ potencialno.
Če nam tako funkcijo uspe najti, smo enačbo rešili, njen prvi integral je enak
\[
  u(x, y) = \int_{(x_0, y_0)}^{(x, y)} \lambda P dx + \lambda Q dy.
\]
Ko iščemo $\lambda$, si lahko pomagamo z enakostjo $\partial_x (\lambda Q) =
\partial_y (\lambda P)$.

Včasih dobimo enačbo oblike $P dx + Q dy = 0$.
Če velja $P_y = Q_x$, je enačba \pojem{eksaktna}, s prvih integralom $u$, za
katerega velja $u_x = P, u_y = Q$.
Rešitev je potem implicitno podana z nivojnico $u = C$.
Če ni eksaktna, poskušamo najti integrirujoči množitelj.

\section{Parametrično reševanje}%

Če imamo implicitno podano NDE prvega reda $F(x, y, y') = 0$, lahko nanjo
gledamo kot na ploskev v $\R^3$ s substitucijo $p = y'$.
Potem lahko enačbo obravnavamo na tri načine:
\begin{itemize}
\item Če $y$ ne nastopa v $F$, parametriziramo nastalo krivuljo z $x = x(t)$ in
  $p = p(t)$.
  Parametrizacijo odvajamo in uporabimo dejstvo $dy = p dx$.
  Iz tega dobimo krivuljo $x \mapsto (x, y(x))$, ki je rešitev enačbe.
\item Če $x$ ne nastopa v $F$, spet parametriziramo $y = y(t)$ in $p = p(t)$,
  ter uporabimo $dy = p dx$.
  Iz tega lahko izračunaš krivuljo $t \mapsto (x(t), y(t))$, ki je rešitev DE\@.
\item Če lahko enega od $x, y, p$ eksplicitno izraziš, ta predpis direktno
  vstaviš v $dy = p dx$, in rešitev izraziš parametrično.
\end{itemize}

\section{Tok}

Naj bo dano vektorsko polje $V$.
\pojem{Tokovnica} tega polja je taka krivulja $\gamma$, da velja
$\dot{\gamma}(t) = V(\gamma(t))$.
\pojem{Tok} je preslikava $(x, t) \mapsto \phi_t(x) \in \R^n$, da je $t \mapsto
\phi_t(x)$ tokovnica za vsak $x$, in $\phi_0(x) = x$ za vse $x$.

Če imamo dano polje v polarnih koordinatah $(\dot{r}, \dot{\varphi}) = W(r,
\varphi)$, ga lahko transformiramo v kartezične koordinate z
\[
  V = (\dot{x}, \dot{y}) = D_\psi (\eta(t)) (\dot{r}, \dot{\varphi}),
\]
kjer je $\eta(t) = (r(t), \varphi(t))$ in $\psi(r, \varphi) = (x, y)$.

\section{Enačbe drugega reda}

Naj bo dan sistem
\begin{gather*}
  \dot{\vec{q}} = \vec{G}(\vec{q}, \vec{p}), \\
  \dot{\vec{p}} = \vec{F}(\vec{q}, \vec{p}).
\end{gather*}
Če obstaja taka funkcija $H: M \subseteq \R^{2n} \to \R$, da je
\begin{gather*}
  \partial_{\vec{q}} H = -\vec{F}, \\
  \partial_{\vec{p}} H = \vec{G},
\end{gather*}
potem pravimo, da je sistem \pojem{hamiltonski}, funkciji $H$ pa pravimo
\pojem{Hamiltonian}.
Tedaj je $H$ tudi prvi integral enačbe.
Da je sistem hamiltonski, mora biti vektorsko polje $(-\vec{F}, \vec{G})$
potencialno.

\section{Cauchyjeve naloge}

\pojem{Cauchyjeva naloga} ali \pojem{začetni problem} je sistem enačb
\begin{gather*}
  y' = f(x, y), \\
  y(x_0) = y_0.
\end{gather*}

\begin{izrek}[Eksistenčni izrek]
  Dana je Cauchyjeva naloga $y' = f(x, y)$, $y(x_0) = y_0$.
  Če je $f: C_{a,b} \to \R$ zvezna in Lipscitzova v spremenljivki $y$, kjer je
  \[
	C_{a,b} = [x_0 - a, x_0 + a] \times [y_0 - b, y_0 + b],
  \]
  obstaja enolična rešitev $\tilde{y}$, definirana na $(x_0 - \alpha, x_0 +
  \alpha)$, kjer je
  \begin{gather*}
	\alpha = \min\{a, \frac{b}{M}, \inv{L}\}, \\
	M = \max_{C_{a,b}} \abs{f(x, y)}.
  \end{gather*}
  Če je $f \in \zvodvi{k}$, je tok tudi $\zvodvi{k}$.
\end{izrek}

Lipscitzov pogoj lahko zamenjamo s pogojem, da je $f$ zvezno odvedljiva.
Izrek je pomemben, ker poleg obstoja zagotavlja tudi enoličnost rešitve.
Če moramo dokazati, da je rešitev tudi globalno enolična, lahko uporabo izreka
verižimo od začetne točke; če dobimo, da je rešitev enolična na razdalji
$\alpha$ od začetne točke, se postavimo v točko $\pol \alpha$, tam uporabimo
izrek, in (pod predpostavko, da se maksimalna razdalja ne zmanjša) dobimo obstoj
in enoličnost do $\frac{3}{2} \alpha$.
Postopek lahko potem nadaljujemo.
Obstoj globalne rešitve pa lahko pokažemo tudi z naslednjo lemo:

\begin{lema}
  Naj rešitev Cauchyjeve naloge obstaja na nekem intervalu $(\sigma, \omega)$ za
  $\sigma, \omega \in [-\infty, \infty]$.
  \begin{itemize}
  \item Če je $\omega < \infty$, potem
	\[
	  \lim_{x \to \omega} \abs{y(x)} = \infty.
	\]
  \item Če je $\sigma > \infty$, potem
	\[
	  \lim_{x \to \sigma} \abs{y(x)} = \infty.
	\]
  \end{itemize}
\end{lema}

\section{Linearni sistemi NDE}

Rešitev sistema $\dot{x} = A x + b$ je vsota homogene in partikularne rešitve.
Za homogen sistem $\dot{x} = Ax$ poiščemo fundamentalno rešitev
\[
  \phi = e^{A t} = P e^{Jt} P^{-1},
\]
kjer je $A=PJP^{-1}$ Jordanov razcep.
Za dobljeno matriko velja $\dot{\phi} = A \phi$ in $\phi(0) = I$, njeni stolpci
podajajo bazo rešitev sistema.
Če je $x(0) = x_0$ začetni pogoj, je $\phi x_0$ rešitev začetnega problema.

Partikularno rešitev poiščemo z nastavkom
\[
  x_p(t) = \int_0^t \phi(t) \phi(s)^{-1} b(s) ds
  = P e^{Jt} \int_0^t e^{-Js} P^{-1} b(s) ds.
\]

\section{Linearne NDE višjega reda}%
\label{sec:linearne-nde-visjega-reda}

Za homogeno enačbo
\[
  a_0 y^{(n)} + a_1 y^{(n-1)} + \ldots + a_n y = 0
\]
zapišemo karakteristični polinom
\[
  f(\lambda) = a_0 \lambda^n + a_1 \lambda^{n-1} + \ldots + a_n = 0.
\]
Rešitve enačbe so vsote funkcij oblike
\[
  p(x) e^{\lambda_i x},
\]
kjer je ničla $\lambda_i \in \R$ in $\deg p = \deg_f \lambda_i - 1$, ter
funkcij oblike
\begin{align*}
  q_1(x) e^{\Re \lambda_i \cdot x} \cos(\Im \lambda_i \cdot x)
  && q_2(x) e^{\Re \lambda_i \cdot x} \sin(\Im \lambda_i \cdot x)
\end{align*}
za $\deg q_1 = \deg q_2 = \deg_f \lambda - 1$ in kompleksno ničlo $\lambda_i$.

Če je desna stran oblike $p(x) e^{\mu x}$, je postopek reševanja opisan v
razdelku~\ref{sec:posebne-nde-visjega-reda}, za splošno desno stran pa vzamemo
linearno neodvisne rešitve homogenega dela $y_1, \ldots, y_n$, in poiščemo še
partikularno rešitev z variacijo konstante
\[
  y_p = C_1(x) y_1 + \ldots + C_n(x) y_n.
\]

\section{Variacijski račun}%
\label{sec:variacijski-racun}

Dana je $L: \R^3 \to \R \in \zvodvi{2}$ ter funkcional
\[
  \mathcal{L}(y) = \int_a^b L(x, y, y') dx
\]
na prostoru vseh $\zvodvi{2}$ funkcij $y: [a, b] \to \R$, za katere veljata
robna pogoja $y(a) = A$ in $y(b) = B$.
Iščemo ekstreme tega funkcionala.
Le-ti zadoščajo Euler-Lagrangeovi enačbi
\[
  L_y(x, y, y') - \partial_x L_{y'}(x, y, y') = 0,
\]
kjer si pri odvajanju po $y$ in $y'$ (indeksa) mislimo, da sta to neodvisni
spremenljivki, pri odvajanju po $x$ ($\partial_x$) pa upoštevamo $y = y(x)$.
Ko enačbo rešiš, dobiš predpis za $y$ z dvema konstantama; konstanti določiš
tako, da predpis vstaviš v robna pogoja.
Če kakšnega robnega pogoja ne poznaš, ga nadomestiš z
\[
  L_{y'}(x_0, y(x_0), y'(x_0)) = 0,
\]
kjer je $x_0$ tisti izmed $a, b$, za katerega ne poznaš začetne vrednosti.

Poznamo tudi dva posebna primera:
\begin{itemize}
\item Če je $L = L(x, y')$, velja $Ly' = \text{konst.}$.
  V tem primeru ne potrebujemo vstavljati v Euler-Lagrangeovo enačbo.
\item Če je $L = L(y, y')$, uporabimo Bertranijevo identiteto $L - y' L_{y'} =
  \text{konst.}$.
  V tem primeru moraš preveriti, da je dobljena rešitev dejansko rešitev
  Euler-Lagrangeove enačbe.
\end{itemize}

Podobno lahko delamo tudi v več dimenzijah; če je $y = (y_1, \ldots, y_n)$
iskana preslikava, preprosto uporabimo Euler-Lagrangeovo za vsako komponento
posebej;
\[
  L_{y_i} - \partial_x (L_{y_i'}) = 0.
\]

Če je $L$ odvisen tudi od odvodov višjih redov, uporabiš Euler-Poissonovo enačbo
\[
  \sum_{k=0}^n (-1)^k (\frac{d}{dx})^k L_{y^{(k)}} = 0.
\]

\section{Pogosti postopki reševanja}%

\subsection{Enačba z ločljivima spremenljivkama}%
\label{sec:locljive-spremenljivke}

Če imamo enačbo oblike
\[
  \dot{x} = f(t) g(x),
\]
jo predelamo v
\[
  \frac{\dot{x}}{g(x)} = f(t),
\]
in integriramo obe strani.
Če je $H$ primitivna funkcija $\nicefrac{1}{g}$, in $F$ primitivna funkcija $f$,
velja
\[
  x(t) = H^{-1}(F(t) + C)
\]
za poljubno konstanto $C \in \R$.

\subsection{Enačba s homogeno desno stranjo}%
\label{sec:homogena-desna-stran}

Če je $\dot{x} = f(t, x)$, kjer velja $f(t, x) = f(\lambda t, \lambda x)$ za vse
$\lambda \in \R\brez\{0\}$, vpeljemo spremenljivko
\begin{align*}
  v = \frac{x}{t} && \dot{x} = t \dot{v} + v.
\end{align*}
Ker velja $\dot{x} = f(1, v)$, dobimo enačbo
\[
  \dot{v} = \inv{t} (f(1, v) - v),
\]
to je enačba z ločljivima spremenljivkama, nadaljuješ kot v
razdelku~\ref{sec:locljive-spremenljivke}.

\subsection{Enačbe višjega reda}%
\label{sec:enacbe-visjega-reda}

Dana je enačba $F(x, y, y', \ldots, y^{(n)}) = 0$.
Imamo tri pristope:
\begin{itemize}
\item Če $F$ ni odvisna od $y$, znižamo red z $z = y'$.
\item Če $F$ ni odvisna od $x$, znižamo red z $z(x) = y'(y)$; pri tem velja
  $\partial_y z = \nicefrac{y''}{y'}$.
\item Če je $F(x, \lambda y, \ldots, \lambda y^{(n)}) = \lambda^k F(x, y,
  \ldots, y^{(n)})$, znižamo red z $z = \nicefrac{y'}{y}$.
\end{itemize}

\subsection{Linearna NDE prvega reda}%
\label{sec:linearna-prvega-reda}

Linearna NDE prvega reda je enačba oblike
\[
  y' = f(x) y + g(x).
\]
Rešiš jo tako, da prvo rešiš homogeno enačbo
\[
  y' = f(x) y,
\]
ki je enačba z ločljivimi spremenljivkami, rešitev poiščeš kot v
razdelku~\ref{sec:locljive-spremenljivke}.
Vedno dobiš
\[
  y_h = C \exp\left( \int_a^x f(\xi) d\xi \right).
\]
Splošna rešitev bo enaka vsoti homogene in partikularne rešitve.
Slednjo poiščemo z nastavkom
\[
  y_p = C(x) \exp\left( \int_a^x f(\xi) d\xi \right),
\]
postopku pravimo \pojem{variacija konstante}.
Če nastavek vstavimo v originalno diferencialno enačbo, dobimo
\[
  C'(x) \exp\left(\int_a^x f(\xi) d\xi \right) = g(x).
\]
Ko to integriramo, dobimo preslikavo $C(x)$, in rešitev
\[
  y = C(x) \exp\left( \int_a^x f(\xi) d\xi \right)
  + D \exp\left( \int_a^x f(\xi) d\xi \right),
\]
kjer je $D \in \R$ poljuben parameter (ki pride od homogene rešitve).

\subsection{Bernoullijeva enačba}%
\label{sec:bernoullijeva-enaba}

Bernoullijeva enačba je enačba oblike
\[
  p(x) y' + q(x) y = r(x) y^\alpha(x)
\]
za nek $\alpha \in \R$.
Če je $\alpha = 0$ ali $\alpha = 1$, sistem rešimo po
razdelku~\ref{sec:linearna-prvega-reda}.
V nasprotnem primeru vpeljemo
\begin{align*}
  z(x) = {(y(x))}^{1 - \alpha}, && z'(x) = (1 - \alpha) y^{-\alpha}(x) y'(x).
\end{align*}
Če začetno enačbo delimo z $y^{-\alpha}$, dobimo
\[
  p y' y^{-\alpha} + q y^{1 - \alpha} = r,
\]
oziroma
\[
  z' + \frac{q}{p} (1 - \alpha) z = \frac{r}{p}(1 - \alpha),
\]
kar je nehomogena NDE prvega reda, ki jo rešimo po
razdelku~\ref{sec:linearna-prvega-reda}.

\subsection{Riccatijeva enačba}%
\label{sec:riccatijeva-enacba}

Riccatijeva enačba je enačba oblike
\[
  y' = a(x) y^2 + b(x) y + c(x),
\]
ki je v splošnem ne znamo rešiti.
Če uganemo neko partikularno rešitev $y_p$, lahko uporabimo nastavek $y = y_p +
z$, ki nam enačbo reducira na
\[
  z' = (2 a y_p + b) z + a z^2,
\]
kar je Bernoullijeva enačba za $p = 1$, $q = (-2a y_p + b)$, $r = a$ in $\alpha
= 2$.
Rešimo jo po postopku v razdelku~\ref{sec:bernoullijeva-enaba}.

\subsection{Clairontova enačba}%
\label{sec:clairontova-enacba}

Clairontova enačba je enačba oblike
\[
  y = xy' + \psi(y').
\]
Rešujemo jo parametrično, torej zapišemo $p = y'$ in uporabimo $dy = p dx$.
Ko dobimo splošno rešitev $y = Cx + \psi(C)$, zapišemo $G(x, y, C) = y - Cx -
\psi(C)$, izračunamo še singularno rešitev z enačbama $G = 0$ in $\partial_C G =
0$.

\subsection{Lagrangeova enačba}%
\label{sec:lagrangeova-enacba}

Lagrangeova enačba je enačba oblike
\[
  y = x \phi(y') + \psi(y').
\]
Rešujemo jo parametrično, torej zapišemo $p = y'$ in uporabimo $dy = p dx$.
S tem se enačba prevede na linearno, ki jo rešimo kot v
razdelku~\ref{sec:linearna-prvega-reda}.

\subsection{Eulerjeva enačba}%
\label{sec:eulerjeva-enacba}

Za reševanje Eulerjeve enačbe
\[
  a_0 x^n y^{(n)} + a_1 x^{n-1} y^{(n-1)} + \ldots + a_n y = 0
\]
uporabimo nastavek $y(x) = x^\lambda$.
V postopku poiščemo ničle polinoma
\[
  q(\lambda) = a_0 \lambda (\lambda-1) \ldots (\lambda - n + 1)
  + a_1 \lambda (\lambda-1) \ldots (\lambda - n + 2)
  + \ldots
  + a_{n-1} \lambda + a_n
\]
Rešitev je vsota funkcij
\[
  p(\ln x) x^\lambda
\]
za realno ničlo $\lambda$ in $\deg p = \deg_f \lambda - 1$.
Za kompleksen $\lambda$ pa
\begin{align*}
  q_1(\ln x) x^{\Re \lambda} \cos(\Im \lambda \cdot x)
  && q_2(\ln x) x^{\Re \lambda} \sin(\Im \lambda \cdot x).
\end{align*}

\subsection{Posebne nehomogene linearne NDE višjega reda}%
\label{sec:posebne-nde-visjega-reda}

Za enačbo oblike
\[
  a_0 y^{(n)} + \ldots + a_{n-1} y' + a_n y = p(x) e^{\mu x}
\]
je rešitev oblike $y_h + y_p$.
Homogen del poiščemo kot v razdelku~\ref{sec:linearne-nde-visjega-reda},
partikularen del pa je oblike
\[
  y_p = e^{\mu x} q(x) x^k,
\]
kjer je $\deg q = \deg p$ in $k = \deg_f \mu$ za polinom $f$ iz
razdelka~\ref{sec:linearne-nde-visjega-reda}.

Podobno rešimo tudi enačbo oblike
\[
  a_0 x^n y^{(n)} + \ldots + a_{n-1} x y' + a_n y = x^\mu p(\ln x)
\]
kot $y = y_h + y_p$, kjer homogen del poiščemo kot v
razdelku~\ref{sec:eulerjeva-enacba}, partikularen pa z
\[
  y_p = x^\mu r(\ln x) (\ln x)^k
\]
za $\deg r = \deg p$ in $k = \deg_q \mu$.
Tu smo se sklicali na polinom $q$ iz razdelka~\ref{sec:eulerjeva-enacba}.

\section{Razno}%

\subsection{Iskanje ortogonalnih krivulj}

Če imamo dano krivuljo $y(x)$, in iščemo nanjo pravokotne krivulje, prvo
odvajamo predpis za to krivuljo, da dobimo enačbo $y = f(x, y')$.
V tej enačbi nato zamenjamo pojavitve $y'$ z $\nicefrac{-1}{y'}$, da dobimo
enačbo za ortogonalne trajektorije.

\subsection{Determinanta Wronskega}

Če imamo homogeno linearno NDE višjega reda
\[
  a_0 y^{(n)} + \ldots + a_{n-1} y' + a_n y = 0,
\]
in zložimo bazne rešitve ter njihove odvode v matriko
\[
  \phi
  =
  \begin{bmatrix}
	y_1 & \cdots & y_n \\
	\vdots & \ddots & \vdots \\
	y_1^{(n-1)} & \cdots & y_n^{(n-1)}
  \end{bmatrix},
\]
lahko zapišemo determinanto Wronskega kot $W(x) = \det \phi(x)$.
Zanjo velja $W'(x) = \operatorname{tr}(A) W(x)$.

V dvodimenzionalnem primeru, ko ima enačba obliko
\[
  y'' + py' + qy = 0,
\]
lahko z determinanto Wronskega poiščemo drugo rešitev, če prvo že imamo.
Če sta $u$ in $v$ dve linearno neodvisni rešitvi, za determinanto $W = uv' -
u'v$ namreč velja $W' + pW = 0$ ter
\[
  v = u \int \frac{W}{u^2} dx.
\]
Če imamo enačbo in eno rešitev, prvo poiščemo $W$ z diferencialno enačbo $W' +
pW = 0$ (če je $p=0$, pride $W=\text{konst.}$, kar je popolnoma sprejemljivo),
nato pa poiščemo drugo rešitev z zgornjim integralom.

% LocalWords:  Wronskega
